\section{Frog Jump with K Distances}
\label{sec:dp-frog-jump-k}

\problemstatement{Same setting as Section~\ref{sec:dp-frog-jump}, except
the frog may jump from index $i$ to any of $i+1, i+2, \ldots, i+k$ (up
to $k$ steps ahead, instead of only $1$ or $2$), at cost
$|\textit{height}[i]-\textit{height}[j]|$ for a jump to index $j$. Find
the minimum total cost to reach index $n-1$.}

\begin{patternbox}
\emph{``Same recurrence shape, but the fixed set of options widens from
$2$ to a parameter $k$''} -- recognizing this is the entire content of
this section: Section~\ref{sec:dp-frog-jump}'s minimum over $2$ candidates
($i-1$ and $i-2$) becomes a minimum over up to $k$ candidates
($i-1, i-2, \ldots, i-k$), with everything else about the derivation
unchanged.
\end{patternbox}

\subsection*{Deriving the recurrence}

The frog's last jump into index $i$ came from some index $i-j$ for some
$j$ between $1$ and $k$ (as long as $i-j \ge 0$). So:
\[
  \textit{minCost}(i) = \min_{1 \le j \le k,\ i-j \ge 0}
  \Big(\textit{minCost}(i-j) + |\textit{height}[i]-\textit{height}[i-j]|\Big).
\]
Base case: $\textit{minCost}(0)=0$.

\subsection*{Approach 1: Recursion}

\begin{lstlisting}
#include <bits/stdc++.h>
using namespace std;

int frogJumpKRecursive(int i, int k, vector<int>& height) {
    if (i == 0) return 0;

    int best = INT_MAX;
    for (int j = 1; j <= k && i - j >= 0; j++) {
        int cost = frogJumpKRecursive(i - j, k, height) +
                    abs(height[i] - height[i - j]);
        best = min(best, cost);
    }
    return best;
}

int main() {
    vector<int> height = {10, 30, 40, 50, 20};
    int n = height.size(), k = 3;
    cout << frogJumpKRecursive(n - 1, k, height) << endl;
    return 0;
}
\end{lstlisting}

\begin{complexitybox}[Approach 1 Complexity]
\textbf{Time.} $O(k^n)$ in the worst case -- each call now branches
into up to $k$ recursive calls instead of $2$.
\textbf{Space.} $O(n)$ for the recursion stack depth.
\end{complexitybox}

\subsection*{Approach 2: Memoization}

\begin{lstlisting}
#include <bits/stdc++.h>
using namespace std;

int frogJumpKMemo(int i, int k, vector<int>& height, vector<int>& memo) {
    if (i == 0) return 0;
    if (memo[i] != -1) return memo[i];

    int best = INT_MAX;
    for (int j = 1; j <= k && i - j >= 0; j++) {
        int cost = frogJumpKMemo(i - j, k, height, memo) +
                    abs(height[i] - height[i - j]);
        best = min(best, cost);
    }
    return memo[i] = best;
}

int main() {
    vector<int> height = {10, 30, 40, 50, 20};
    int n = height.size(), k = 3;
    vector<int> memo(n, -1);
    cout << frogJumpKMemo(n - 1, k, height, memo) << endl;
    return 0;
}
\end{lstlisting}

\begin{complexitybox}[Approach 2 Complexity]
\textbf{Time.} $O(n \cdot k)$ -- $n$ distinct sub-problems, each doing
$O(k)$ work in its loop.
\textbf{Space.} $O(n)$ memo $+$ $O(n)$ recursion stack.
\end{complexitybox}

\subsection*{Approach 3: Tabulation}

\begin{lstlisting}
#include <bits/stdc++.h>
using namespace std;

int frogJumpKTabulation(vector<int>& height, int k) {
    int n = height.size();
    vector<int> dp(n, 0);

    for (int i = 1; i < n; i++) {
        int best = INT_MAX;
        for (int j = 1; j <= k && i - j >= 0; j++) {
            int cost = dp[i - j] + abs(height[i] - height[i - j]);
            best = min(best, cost);
        }
        dp[i] = best;
    }
    return dp[n - 1];
}

int main() {
    vector<int> height = {10, 30, 40, 50, 20};
    cout << frogJumpKTabulation(height, 3) << endl;
    return 0;
}
\end{lstlisting}

\subsection*{Dry run of the tabulation table}

\dryrun{Take $\textit{height} = [10,30,40,50,20]$, $k=3$.}

\begin{itemize}
  \item $\textit{dp}[0]=0$.
  \item $\textit{dp}[1] = \textit{dp}[0]+|30{-}10| = 20$ (only $j=1$
        possible).
  \item $\textit{dp}[2] = \min(\textit{dp}[1]+|40{-}30|,\
        \textit{dp}[0]+|40{-}10|) = \min(30,30) = 30$.
  \item $\textit{dp}[3] = \min(\textit{dp}[2]+|50{-}40|,\
        \textit{dp}[1]+|50{-}30|,\ \textit{dp}[0]+|50{-}10|) =
        \min(40,40,40) = 40$.
  \item $\textit{dp}[4] = \min(\textit{dp}[3]+|20{-}50|,\
        \textit{dp}[2]+|20{-}40|,\ \textit{dp}[1]+|20{-}30|) =
        \min(40{+}30,\ 30{+}20,\ 20{+}10) = \min(70,50,30) = 30$.
\end{itemize}

\begin{center}
\begin{tabular}{c|c|c|c|c|c}
$i$ & 0 & 1 & 2 & 3 & 4 \\ \hline
$\textit{dp}[i]$ & 0 & 20 & 30 & 40 & 30
\end{tabular}
\end{center}

The answer is $\textit{dp}[4]=30$, achieved by jumping directly from
index $1$ ($\textit{height}=30$) to index $4$ ($\textit{height}=20$), a
jump of size $3 \le k$, costing $\textit{dp}[1] + |20-30| = 20+10=30$.

\begin{complexitybox}[Approach 3 Complexity]
\textbf{Time.} $O(n \cdot k)$.
\textbf{Space.} $O(n)$ table, no recursion stack.
\end{complexitybox}

\subsection*{Approach 4: Space Optimization}

\begin{ideabox}[Why This One Cannot Be Reduced All the Way to $O(1)$]
Unlike Section~\ref{sec:dp-frog-jump}, where only the \emph{two}
immediately preceding entries were ever needed, here computing
$\textit{dp}[i]$ may need any of the previous $k$ entries -- so the
space optimization here reduces storage to a \textbf{sliding window} of
the last $k$ entries (e.g.\ a small circular buffer or deque of size
$k$), rather than collapsing all the way to two scalar variables. When
$k$ is a fixed small constant, this is still $O(k) = O(1)$ relative to
$n$; when $k$ can grow proportionally with $n$, no reduction below
$O(n)$ is possible, since up to $n$ of the entries might genuinely be
needed.
\end{ideabox}

\begin{lstlisting}
#include <bits/stdc++.h>
using namespace std;

int frogJumpKSpaceOptimized(vector<int>& height, int k) {
    int n = height.size();
    vector<int> window(k, 0); // window[i % k] holds dp[i]
    window[0] = 0;

    for (int i = 1; i < n; i++) {
        int best = INT_MAX;
        for (int j = 1; j <= k && i - j >= 0; j++) {
            int cost = window[(i - j) % k] + abs(height[i] - height[i - j]);
            best = min(best, cost);
        }
        window[i % k] = best;
    }
    return window[(n - 1) % k];
}

int main() {
    vector<int> height = {10, 30, 40, 50, 20};
    cout << frogJumpKSpaceOptimized(height, 3) << endl;
    return 0;
}
\end{lstlisting}

\begin{complexitybox}[Approach 4 Complexity]
\textbf{Time.} $O(n \cdot k)$, unchanged.
\textbf{Space.} $O(k)$ for the sliding window, instead of $O(n)$ for
the full table.
\end{complexitybox}

\begin{takeawaybox}
Generalizing a recurrence's fixed-size option set from a small constant
(here, $2$, in Frog Jump) to a parameter $k$ widens every stage's loop
from $2$ iterations to $k$, multiplying the time complexity by $k$ and
capping how far the space optimization can shrink: only the last $k$
table entries, not just the last $2$, must remain available, so the
``two rolling variables'' trick of earlier sections becomes a
``sliding window of size $k$'' trick here.
\end{takeawaybox}
